\section{Line A: Physiological Information in Facial Video}
Three rPPG estimators were applied to the video of the four pre-registered UBFC-Phys participants (s1--s4; 12 clips), and their pulse estimates were compared between stress and rest with one-sided Mann--Whitney $U$ tests.

\begin{table}[htbp]
\centering
\begin{tabular}{@{}lccl@{}}
\toprule
\textbf{rPPG estimator} & \textbf{$U$} & \textbf{$p$-value} & \textbf{Status} \\
\midrule
POS ($I = 0.5$) & 1050.0 & 0.0314 & Primary pre-registered test \\
CHROM & 1318.0 & 0.0001 & Exploratory (post hoc) \\
PBV & 1113.0 & 0.0103 & Exploratory (post hoc) \\
\bottomrule
\end{tabular}
\caption{Tests for physiological information recoverable from facial video (UBFC-Phys, participants s1--s4).}
\label{tab:leakage}
\end{table}

The pre-registered primary test (POS, $U = 1050.0$, $p = 0.0314$) rejects, at the pre-registered threshold $p \le 0.05$, the null hypothesis that video compression removes all recoverable physiological signal. The two exploratory estimators point in the same direction but are not treated as independent confirmation. In a per-clip analysis, the CHROM estimate for participant s2 during T1 correlated with the wrist BVP reference at $r = 0.849$, with matching cardiac frequency and a spectral signal-to-noise ratio of 53.5, indicating that the recovered signal is physiological rather than a statistical artifact. The facial video therefore carries genuine information about the physiological state that defines the stress label, so a visual branch trained on real video could legitimately use it.

The dataset authors recommend excluding some clips from rPPG evaluation because of video quality; within this cohort, these are s1 T2, s3 T1, and s4 T2. The pre-registered analysis included these clips; an exploratory analysis excluding them will be reported in the final version.

\paragraph{Sensitivity analysis.} A post hoc audit found that two clips (s4 T2 and s4 T3) in the $I = 1.0$ ablation arm had face crops about 37\% smaller than the pre-registered crop-size rule required, because of a detection anomaly. Excluding them from that arm gives $U = 832.0$, $p = 0.0857$, above the threshold, so the $I = 1.0$ result ($p = 0.0437$) is fragile and is not cited as independent evidence. The primary $I = 0.5$ result and the exploratory CHROM and PBV results are unaffected by this anomaly.

\section{Line B: Invariance to the Synthetic Scar}
\subsection{EQUITAS-RCMF across Regimes}
Table~\ref{tab:regimes} reports EQUITAS-RCMF on the 496 test samples from three held-out WESAD participants, in each regime and mode.

\begin{table}[htbp]
\centering
\begin{tabular}{@{}llcccc@{}}
\toprule
\textbf{Mode} & \textbf{Regime} & \textbf{Accuracy} & \textbf{DP gap} & \textbf{EO gap} & \textbf{CF gap} \\
\midrule
Autonomous & Biased ($\rho = 0.85$) & 0.617 & 0.161 & 0.053 & 0.00064 \\
Autonomous & Neutral ($\rho = 0.50$) & 0.617 & 0.048 & 0.065 & 0.00064 \\
Autonomous & Inverted ($\rho = 0.15$) & 0.617 & 0.177 & 0.042 & 0.00064 \\
\midrule
Privileged & Biased ($\rho = 0.85$) & 0.625 & 0.169 & 0.026 & 0.0024 \\
Privileged & Neutral ($\rho = 0.50$) & 0.627 & 0.060 & 0.097 & 0.0024 \\
Privileged & Inverted ($\rho = 0.15$) & 0.627 & 0.190 & 0.039 & 0.0024 \\
\bottomrule
\end{tabular}
\caption{EQUITAS-RCMF on the Line~B test set (seed 42; 496 samples from participants S5, S10, and S11). Autonomous mode is the deployed, mask-free configuration. Source: \texttt{equitas\_rcmf\_master\_benchmark\_report.json}.}
\label{tab:regimes}
\end{table}

In autonomous mode, accuracy is 0.617 in all three regimes, and removing the scar changes the predicted probability of stress by 0.00064 on average. Reversing the association between the scar and the label between training and test therefore leaves the model's performance unchanged, which is the behaviour expected of a model that does not use the scar. Privileged mode gives similar accuracy (0.625--0.627) with a larger but still small counterfactual gap (0.0024).

The DP gap rises from 0.048 in the neutral regime to 0.161 and 0.177 in the biased and inverted regimes. This is expected rather than a sign of unfairness. When the scar is assigned according to the label, scarred and unscarred inputs have different proportions of stress windows, so even a model that ignores the scar produces different rates of positive predictions for the two groups. The appropriate criterion in this setting is equalized odds, which compares error rates within each class~\cite{anthis2023}; the EO gap stays at or below 0.065 in autonomous mode across all three regimes.

The neutral regime also corresponds to the setting analysed by Anthis and Veitch~\cite{anthis2023}: a model trained where a spurious attribute is associated with the label ($\rho = 0.85$) and evaluated where it is not. A model that relies on the scar would lose accuracy in this shift, whereas EQUITAS-RCMF does not.

\subsection{Comparison Models}
\paragraph{Model A (Naive ERM).} Results for the naive ERM model under the three regimes will be reported in the final version.
% TODO(authors): If thesis_production_benchmark_report.json contains the "baseline"
% (naive ERM) results for each regime on this benchmark, add them here as a table next to
% the regimes table. The earlier validation accuracies (79.94-80.98%) must not be
% used unless their reports' csv_path is multimodal_publishable.csv.

\paragraph{Model B (camera-off).} A physiology-only model trained and evaluated on this benchmark's split will be reported in the final version. The existing camera-off result (72.15\% validation accuracy) was obtained on an earlier dataset and is not comparable with Table~\ref{tab:regimes}.

\paragraph{Model C (CGF).} A model trained with a counterfactual penalty but without the subspace decomposition reached 50.20\% accuracy, close to chance, with a CF gap of 0.0019. A counterfactual penalty alone therefore reduced the counterfactual gap but did not produce a useful classifier.
% TODO(authors): State the regime and mode of the Model C figures, and reconcile its DP
% and EO gaps (the report shows 0.0040 and 0.0346) before adding them.

\subsection{Face-Attribute Audit}
Table~\ref{tab:demographic_audit} reports EQUITAS-RCMF in autonomous mode on the neutral regime, grouped by the estimated gender and age of the face images.

\begin{table}[htbp]
\centering
\begin{tabular}{@{}lccccc@{}}
\toprule
\textbf{Estimated attribute} & \textbf{Faces} & \textbf{Accuracy} & \textbf{DP gap} & \textbf{EO gap} & \textbf{CF gap} \\
\midrule
Female & 44 & 62.50\% & 0.0404 & 0.0573 & 0.0007 \\
Male & 18 & 59.72\% & 0.0589 & 0.1212 & 0.0004 \\
Age 18--30 & 28 & 60.27\% & 0.0187 & 0.0581 & 0.0007 \\
Age 30--45 & 21 & 64.29\% & 0.0034 & 0.0621 & 0.0006 \\
Age 45--65 & 13 & 60.58\% & 0.1919 & 0.2308 & 0.0006 \\
\bottomrule
\end{tabular}
\caption{EQUITAS-RCMF by estimated face attributes (seed 42, autonomous mode, $\rho = 0.50$). Gender and age are estimated from the FFHQ images by InsightFace and describe the face images, not the WESAD participants; all stress labels in the test set come from three participants. Source: \texttt{equitas\_rcmf\_master\_benchmark\_report.json}.}
\label{tab:demographic_audit}
\end{table}

Across these groups, the counterfactual gap stays between 0.0004 and 0.0007. The DP and EO gaps vary more, reaching 0.1919 and 0.2308 for the 45--65 group, which contains only 13 faces. Because the attributes are automated estimates and the groups are small, these figures indicate how the model's visual branch behaves on different kinds of face images rather than how it treats different groups of people.

\section{Explainability Analysis}
Integrated Gradients~\cite{sundararajan2017} attributions were computed for the ERM baseline and for EQUITAS-RCMF. For the ERM model, attribution was concentrated on the brow-line scar; for EQUITAS-RCMF, attribution on the scar was negligible. These maps are consistent with the counterfactual results, but they are treated as supporting evidence only: Integrated Gradients maps can remain visually plausible even when they no longer reflect what a model has learned, and thin, high-contrast structures such as a scar are the kind of edge that input-dominated attributions tend to highlight~\cite{adebayo2018}.
% TODO(authors): Confirm which images the attribution figure uses. If they are Line B
% (FFHQ) images, do not describe the attribution as following a pulse signal, since
% still images contain none.

\section{Perturbation-Based Sanity Check}
A sanity-check protocol was implemented (\texttt{src/evaluation/sham\_edit\_probe.py}) that applies a $30 \times 30$-pixel non-semantic perturbation at a fixed location of clean test images and measures the change in prediction and attribution. This protocol complements, but differs from, the parameter- and label-randomization tests of Adebayo et al.~\cite{adebayo2018}. Its results, and those of a parameter-randomization test, will be reported in the final version.

\section{Edge Deployment Benchmark}
\label{sec:deployment}
For deployment, the model was compiled with \texttt{torch.jit.trace} in autonomous mode, with no mask provided. Tracing records only the operations executed for this input, so the mask input and the privileged focus computation are absent from the compiled graph (\texttt{equitas\_rcmf\_edge\_compiled.pt}). The deployed model therefore does not require or accept a scar mask. It still processes the full face image, including the scar region, through the backbone and both subspace projections; its invariance to the scar depends on what the network learned during training, not on the absence of operations in the graph.

Latency was measured on an AMD Ryzen~5 7500F desktop CPU after 50 warm-up iterations, over 1,000 iterations:
\begin{table}[htbp]
\centering
\begin{tabular}{@{}lc@{}}
\toprule
\textbf{Metric} & \textbf{Value} \\
\midrule
Mean latency & 0.90 ms/frame \\
95th-percentile latency & 1.30 ms/frame \\
Throughput & 1112.7 frames/s \\
Real-time budget (30 frames/s) & 33.3 ms/frame \\
Margin relative to budget & 37.0$\times$ \\
\bottomrule
\end{tabular}
\caption{Inference latency of the compiled EQUITAS-RCMF model.}
\label{tab:latency}
\end{table}
These measurements show that the model is computationally light; they do not establish its performance on a specific embedded device.
% TODO(authors): State the batch size, number of threads, and whether the latency
% includes preprocessing (face detection, cropping) or only the network.

\section{Limitations}
\label{sec:limitations}
Several limitations bound these results.
\begin{itemize}
    \item \textbf{Small cohorts.} Line~A uses four participants and 12 clips, so significance tests were carried out at the clip level and equivalence testing was withdrawn. In Line~B, all test labels come from three participants, so estimates of variation across people, including bootstrap intervals clustered by participant, are not informative.
    \item \textbf{Missing physiology-only comparison.} Because the faces carry no stress information, the best achievable accuracy on this benchmark is that of a model using physiology alone. Until a physiology-only model is evaluated on the same split, it is not possible to say whether EQUITAS-RCMF's invariance comes at a cost in accuracy.
    \item \textbf{Focus-scale difference.} The privileged focus used in training and the autonomous focus used at inference lie on different scales, so the deployed gate may penalize scar-focused visual features less strongly than the trained gate. The autonomous-mode results reflect the deployed behaviour, but the gate's training signal and its deployed input are not matched.
    \item \textbf{Estimated attributes.} The face-attribute audit uses gender and age estimated by InsightFace from the face images, which are approximate and do not describe the participants whose physiology defines the labels.
    \item \textbf{Separate lines of evidence.} Line~A shows that real facial video carries physiological information, while Line~B uses face images that carry none, so that the scar is the only label-related visual cue. The benchmark therefore tests whether a model can ignore a spurious attribute; it cannot test whether the model makes good use of legitimate visual evidence, which would require real video with controlled attributes.
    \item \textbf{Single seed.} The Line~B results are for one training seed.
    % TODO(authors): Update if other EQUITAS-RCMF seeds are evaluated.
\end{itemize}
These limitations call for caution in extrapolating the results to larger cohorts and uncontrolled deployment conditions.
